\documentclass[12pt]{article}
\usepackage{hyperref}

\usepackage{graphicx}

\usepackage{mathtools}
\DeclarePairedDelimiter\ceil{\lceil}{\rceil}
\DeclarePairedDelimiter\floor{\lfloor}{\rfloor}


\usepackage{amsmath}
\usepackage{amsfonts}
\usepackage{amssymb}

\usepackage{listings}
\lstset{language=C++,
                basicstyle=\ttfamily,
                keywordstyle=\color{blue}\ttfamily,
                stringstyle=\color{red}\ttfamily,
                commentstyle=\color{green}\ttfamily,
                morecomment=[l][\color{magenta}]{\#}
}

\usepackage{breqn}
\usepackage[]{algorithm2e}
\usepackage{physics}

\newcommand\fnurl[2]{%
  \href{#2}{#1}\footnote{\url{#2}}%
}

\newcommand{\Four}{\mathcal{F}}
\renewcommand{\(}{\left(}
\renewcommand{\)}{\right)}
%% \renewcommand{\{}{\left{}
%% \renewcommand{\}}{\right}}


\title{Real-Complex FFT Notes for rocFFT}
\author{The rocFFT Team}



\def\no#1{\nobreak{#1}} % stop equation breaking in dmath

\begin{document}
\maketitle

\section{Introduction}

This document contains notes for computing real/complex FFTs
for \texttt{rocFFT}.


\subsection{Properties of real/complex DFTs}

The output of a real-to-complex Fourier transform exhibits Hermitian
symmetry, which is to say
\begin{dmath}
  F_k = F^*_{-k}
\end{dmath}
where ${}^*$ denotes complex conjugation and $k$ is a (possibly
multi-dimensional) wave vector.  Note that $k$ is considered to be a
cyclic in each index.

It is easy to see that the origin is real-valued, since
$F_0=F^*_{-0}=F^*_{0}\implies F_0\in\mathbb{R}$.

In adidtion to restrictions on the origin, we may have verious
restrictions on the Nyquist frequency (or frequencies, in the case of
multi-dimensional transforms).

\subsubsection{1D Hermitian symmetry}
For 1D transforms of length $N_x$, we have two cases:
\begin{itemize}
\item $N_x\mid 2$.  Then $F_{N_x/2} = F^*_{-N_x/2} = F^*_{N_x/2}$, so
  $F_{N_x/2}\in\mathbb{R}$.
\item $N_x\nmid 2$.  Then the highest fequency is $\frac{N_x - 1}{2}$,
  and is conjugate to
  $\frac{-N_x+1}{2}=\frac{N_x+1}{2}=neq\frac{N_x-1}{2}$, so there is no reality
  condition for the Nyquist frequency in this case.
\end{itemize}

\subsubsection{2D Hermitian symmetry}
For 2D transforms of length $N_x\times{}N_y$, we store indices $(i,j)$ with
$i=0,\dots{}N_x-1$ and $j=0,\dots{}\floor{N_y/2}$.  We have:
\begin{itemize}
\item $F_{0,0}\in\mathbb{R}$.
\item If $N_x\mid 2$, then $F_{\frac{N_x}{2},0}\in\mathbb{R}$
\item If $N_y\mid 2$, then $F_{0,\frac{N_y}{2}}\in\mathbb{R}$
\item If $N_x\mid 2$ and $N_y\mid 2$, then
  $F_{\frac{N_x}{2},\frac{N_y}{2}}\in\mathbb{R}$
\end{itemize}
Moreover,
\begin{itemize}
\item Since we keep all of the $x$-axis, $F_{i,0}=F^*_{N_x-1,0}$ for
  $i=1,\dots,\floor{N_x/2}-1$.
\item If $N_y\mid 2$, then
  $F_{i,N_y/2}=F^*_{-i,-N_y/2}=F^*_{N_x-i,N_y/2}$, so $F_{i,N_y/2}$
  and $F_{N_x-i,N_y/2}$ are conjugate-pairs for
  $i=1,\dots,\floor{N_x/2}-1$.
\end{itemize}
However, there is no such condition if $N_x\mid 2$, because
$(\frac{N_x}{2},j)\rightarrow(-\frac{N_x}{2},-j)=(\frac{N_x}{2},N_y-j)$,
and, for $j=0,\dots{},\floor{N_y/2}-1$, $N_y-j$ isn't explicitly
stored.

\subsubsection{3D Hermitian symmetry}

For 3D transforms of length $N_x\times{}N_y\times{}N_z$, we store
indices $(i,j,k)$ with $i=0,\dots,N_x-1$, $j=0,\dots,N_y-1$, and
$k=0,\dots\floor{N_z/2}$.

\begin{itemize}
\item $F_{0,0,0} \in \mathbb{R}$ .
\item If $N_x \mid 2$, then $F_{N_x/2,0,0} \in \mathbb{R}$
\item If $N_y \mid 2$, then $F_{0,N_y/2,0} \in \mathbb{R}$
\item If $N_z \mid 2$, then $F_{0,0,N_z/2} \in \mathbb{R}$
\item If $N_x \mid 2$ and $N_y \mid 2$, then $F_{N_x/2,N_y/2,0} \in \mathbb{R}$
\item If $N_x \mid 2$ and $N_z \mid 2$, then $F_{N_x/2,0,N_z/2} \in \mathbb{R}$
\item If $N_y \mid 2$ and $N_z \mid 2$, then $F_{0,N_y/2,N_z/2} \in \mathbb{R}$
\item If $N_x \mid 2$, $N_y \mid 2$, and $N_z \mid 2$, then
  $F_{N_x/2,N_y/2,N_z/2} \in \mathbb{R}$
\end{itemize}
Moreover, then $k=0$ we have redundant data which must be made
consistent.
\begin{itemize}
\item Since we keep all of the $x$-axis, $F_{i,0,0}=F^*_{N_x-1,0,0}$ for
  $i=1,\dots,\floor{N_x/2}-1$.
\item Now the $x$--$y$ plane: $F_{i,j,0}=F^*_{N_x-i,N_y-j,0}$ for
  $i=1,\dots,N_x-1$, $j=1,\dots,\floor{N_y/2}-1$.
\item Similarly, if $N_z\mid{}2$, then
  \begin{itemize}
  \item $F_{i,0,N_z/2}=F^*_{N_x-1,0,N_z/2}$ for $i=1,\dots,\floor{N_x/2}-1$.
  \item $F_{i,j,N_z/2}=F^*_{N_x-i,N_y-j,N_z/2}$ for $i=1,\dots,N_x-1$,
    $j=1,\dots,\floor{N_y/2}-1$
  \end{itemize}
\end{itemize}


\section{1D Transforms}

\subsection{Using a complex array with zero imaginary part}
\label{embed1d}
The easiest way is to just compute a length-$N$ complex FFT with zero
imaginary part.  Using the split format, one can allocate $N$ real
values and use this buffer to hold the imaginary part of the
split-array complex format.  Setting this imaginary array to be all
zero allows one to compute the real-to-complex transform by just
performing a complex-to-complex transform and amalgamating the output
into the format we want.

\subsection{Packing two $\mathbb{R}$ transforms into one $\mathbb{C}$ transform}
\label{packed1d}

Let $x=\left\{x_\ell\right\}_{\ell=0}^{N-1}$ and
$y=\left\{y_\ell\right\}_{\ell=0}^{N-1}$, be real-valued
sequences. Let $z=x + iy$. By the linearity of the Fourier transform,
$Z=X+iY$. Since $x$ and $y$ are real-valued, we have $X_r^*=X_{N-r}$.
Thus,
\begin{dmath}
  \frac{Z_r + Z_{N-r}^*}{2}
  = \frac{X_r + i Y_r}{2} + \(\frac{X_{N-r} + iY_{N-r}}{2}\)^*
  = \frac{X_r + i Y_r}{2} + \(\frac{X^*_r + iY^*_r}{2}\)^*
  = \frac{X_r + i Y_r}{2} + \frac{X_r - iY_r}{2}
  =  X_r.
\end{dmath}
Similarly,
\begin{dmath}
  \frac{Z_r - Z_{N-r}^*}{2i} =  Y_r.
\end{dmath}
Thus we can compute the Fourier transform of $x$ and $y$ for
$r=1,\dots,\floor*{\frac{N}{2}}+1$ via
\begin{equation}\label{packedrcfft}
  \boxed{X_r = \frac{Z_r + Z_{N-r}^*}{2}},\qquad
  \boxed{Y_r = \frac{Z_r - Z_{N-r}^*}{2i}}.
\end{equation}
For $r$=0, we use the fact that $Z$ is cyclic, so $Z_N=Z_0$, and
\begin{equation}\label{packedrcfft0}
  \boxed{X_0 = \Re{Z_0},\qquad  Y_0 = \Im{Z_0}.}
\end{equation}
Since $X_r^*=X_{N-r}$, we don't have to store all of the data in
equation~\eqref{packedrcfft}, but just for
$r=0,\dots,\floor*{N/2}-1$.  The data flow diagram for this
technique is given in Figure~\ref{packedfftflow}.
\begin{figure}
  \centering
  \includegraphics[width=\textwidth]{packedfft}
  \caption{Data flow diagram for packed real-to-complex transform
    given in equation~\eqref{packedrcfft}.}
  \label{packedfftflow}
\end{figure}
The inverse transform $(X,Y)\rightarrow(x,y)$ is performed by
computing $Z$ from $X$ and $Y$ via
\begin{align}
  Z_0 &= \Re{X_0} + i \Re{Y_0}\\
  \label{packinvr}
  Z_r &= X_r + iY_r, \qquad r = 1,\dots, \floor*{\frac{N}{2}} \\
  \label{packinvNr}
  Z_{N-r} &= X_{r}^* + iY_{r}^*, \qquad r = 1,\dots, \floor*{\frac{N}{2}}.
\end{align}
Note that, for even $N$, if $r=N/2$, then $N-r=r$, and
equations~\eqref{packinvr} and~\eqref{packinvNr} are consistent but
redundant.  The data flow diagram for an inverse packed FFT is shown
in Figure~\ref{packedifftflow}.
\begin{figure}
  \centering
  \includegraphics[width=\textwidth]{packedifft}
  \caption{Data flow diagram for complex-to-real packed FFT.}
  \label{packedifftflow}
\end{figure}


\subsection{Computing one even-length $\mathbb{R}$ transform}
\label{even1d}

Let $x=\left\{x_\ell\right\}_{r=0}^{N-1}$ be the real-valued input and
suppose $N\mid{}2$. Let $X=\Four_N\(x\)$ be the length-$N$ transform
of $x$.  Let $x_{\rm{e}}=\left\{x_{2\ell}\right\}$ and
$x_{\rm{o}}=\left\{x_{2\ell+1}\right\}$, and
\begin{dmath}
  z_\ell = x_{2\ell} + i x_{2\ell + 1},
  \quad \text {for } \ell \no= 0,\dots, \frac{N}{2} -1
\end{dmath}
and denote $Z=\Four_{N/2}\(z\)$, $E=\Four_{N/2}\(x_{\rm e}\)$, and
$O=\Four_{N/2}\(x_{\rm o}\)$.  From the linearity of the Fourier
transform, $Z=E+iO$.  Since $x_{\rm{e}}$ and $x_{\rm{o}}$ are real,
their transforms are Hermitian-symmetric (ie $E_{N/2-r}=E_r^*$ and
$O_{N/2-r}=O_r^*$), one can easily verify that
\begin{align}
  E_r &= \frac{1}{2}\(Z_r + Z_{\frac{N}{2}-r}^*\), \\
  O_r &= \frac{i}{2}\(Z_{\frac{N}{2}-r}^*-Z_r\).
\end{align}
We can compute a length--$N$ FFT from two length--$N/2$ FFTs via
\begin{dmath}
  X_r
  = \sum_{\ell=0}^{N-1} x_\ell \omega_N^{r\ell}
  = \sum_{\ell=0}^{N/2-1} x_{2\ell} \omega_{N/2}^{r\ell}
  + \omega_N^r  \sum_{\ell=0}^{N/2-1} x_{2\ell+1} \omega_{N/2}^{r\ell}
  = E_r + \omega_N^r O_r,
\end{dmath}
which simplifies to
\begin{dmath}
  X_r
  = \frac{1}{2}\(Z_r + Z_{\frac{N}{2}-r}^*\)
  + \omega_N^r \frac{i}{2}\(Z_{\frac{N}{2}-r}^*-Z_r\),
\end{dmath}
so we can compute $X$ from $Z$ via

\begin{dmath}\label{rcfftpost}
  \boxed{
    X_r
    = Z_r \frac{1 - i\omega_N^r}{2} + Z_{\frac{N}{2}-r}^* \frac{1 + i\omega_N^r}{2}
    \quad\text{ for } r = 1,\dots\,\frac{N}{2}-1.}
\end{dmath}
We can compute $X_{N/2}$ via the fact that $Z_{N/2}=Z_0$, so
\begin{dmath}
  X_{\frac{N}{2}}
  = Z_{\frac{N}{2}} \frac{1 - i\omega_N^{N/2}}{2}
  + Z_{\frac{N}{2}-\frac{N}{2}}^* \frac{1 + i\omega_N^{N/2}}{2}
  = Z_0 \frac{1 + i}{2} +  Z_0^* \frac{1 -i}{2}
%  = \Re{Z_0} - \Im{Z_0}.
\end{dmath}
which gives
\begin{dmath}\label{XN}
  \boxed{X_{\frac{N}{2}} = \Re{Z_0} - \Im{Z_0}.}
\end{dmath}
Similarly, 
\begin{dmath}\label{X0}
  \boxed{X_0 
    %= Z_0 \frac{1 - i\omega_N^0}{2} +  Z_{\frac{N}{2}-0}^* \frac{1 + i\omega_N^0}{2}
  %= Z_0 \frac{1 - i}{2} + Z_0^* \frac{1 + i}{2}
  = \Re{Z_0} + \Im{Z_0}.}
\end{dmath}
\begin{figure}
  \centering
  \includegraphics[width=\textwidth]{rcffteven}
  \caption{Data flow diagram for even-length real-to-complex FFT.}
  \label{rcffteven}
\end{figure}

To compute a Hermitian-symmetric to real-valued FFT, one inverts
equation~\eqref{rcfftpost} to get $Z_r$ from $X_r$.  From
equation~\eqref{rcfftpost}, we have
\begin{eqnarray}
  \label{2xr}
  2 X_r &=& Z_r \(1 - i \omega_N^r \) + Z_{N/2 -r}^* \(1 + i \omega_N^r \)
  \\
  \label{2xN2r}
  2 X_{N/2 -r}^* &=& Z_r \( 1 + i \omega_N^r \) + Z_{N/2 -r}^* \( 1 - i \omega_N^r \)
\end{eqnarray}
Taking equation~\ref{2xr} times $\(1-i\omega_N^r\)$ minus
equation~\ref{2xN2r} times $\(1+i\omega_N^r\)$, we get
\begin{dmath}
  Z_r
  = X_r \frac{1 + i \omega_N^{-r}}{2}
   + X^*_{\frac{N}{2}-r} \frac{1 - i \omega_N^{-r}}{2}
  %Z_r = X_r\frac{1+i\omega_N^{-r}}{2} - X_{N/2-r}^*\frac{1-i\omega_N^{-r}}{2}.
\end{dmath}
From equations~\eqref{XN} and~\eqref{X0}, we also have, using the fact
that $X_{N/2}$ is real-valued, that
\begin{dmath}
  Z_0 = \frac{X_0 + X_{N/2}}{2} + i \frac{X_0 - X_{N/2}}{2}.
\end{dmath}
However, note that these are the normalized transforms.  Removing the
factor of 2, we get
\begin{eqnarray}
  Z_r
  &=& X_r \(1 + i \omega_N^{-r}\)
  + X^*_{\frac{N}{2}-r} \(1 - i \omega_N^{-r}\), \quad r = 1,\dots,N-1 \\
  Z_0 &=& X_0 + X_{N/2} + i \(X_0 - X_{N/2}\)
\end{eqnarray}
which can be reduced to one case:
\begin{dmath}\label{xr}
\boxed{Z_r 
  = X_r \(1 + i \omega_N^{-r}\)
  + X^*_{\frac{N}{2}-r} \(1 - i \omega_N^{-r}\), \quad r = 0,\dots,N-1}.
\end{dmath}
which gives us the pre-processing stage for the complex-to-real FFT
when the problem size is even.  The data flow diagram is given in
Figure~\ref{crffteven}.
\begin{figure}
  \centering
  \includegraphics[width=\textwidth]{crffteven}
  \caption{Data flow diagram for even-length complex-to-real FFT.}
  \label{crffteven}
\end{figure}


\subsubsection{Strides}
The real-to-complex data flow diagram in Figure~\ref{rcffteven} starts
with computing $z$ from $x$.  When the input stride is 1, this is a
simple cast operation.  When the input stride is greater than 1 more
work is required, as the real and complex parts of $z$ are no longer
contiguous in memory.  One might be able to abuse the split array
format\footnote{\href{http://www.fftw.org/fftw3_doc/Interleaved-and-split-arrays.html}{FFTW},
  \href{https://docs.nvidia.com/cuda/cufft/index.html}{cuFFT}}, where
a complex array is stored in separate real and complex buffers.  One
points the real buffer to $x_0$ and the imaginary buffer to $x_1$ and
then adjusts the strides accordingly.  From the rocFFT
documentation\footnote{\href{https://rocfft.readthedocs.io/en/latest/library.html\#transform-and-array-types}{rocFFT:
    Transform and Array types}}, one is led to believe that the
complex planar format supports this already.  Note that, as of
2019-05-27, the complex-to-complex transform does not support
\lstinline{rocfft_array_type_complex_planar}.


\subsection{$N\nmid{}2$ real-complex transforms}

Q: if, for example, $N\nmid{}3$, can we somehow pack 3 FFTs together?

Alternatively, we can probably just put the data into a complex array
with zero imaginary part and then do a the FFT on that, keeping just
the part that we need.

If $N$ is prime, then we need to do Bluestein.  However, to zero-pad
from $N$ to $2N$ for the Bluestein convolution gives us a problem size
that is divisible by 2, which means that we can presumably do both the
insertion into a larger buffer and the even/odd to complex packing at
the same time, at which point we now have a complex convolution
(perhaps).

Q: for convolutions of real input, can we just abuse the multiply and
avoid having to do all of the extra work in the post/pre processing
stage?  I kind of doubt it, but worth checking out.

\section{Multi-dimensional transforms}

Multidimensional transforms apply Hermitian symmetry in the last
dimension.  For 2D transforms, this is the $y$-direction, and for 3D
this is the $z$-direction.

\subsection{Computation via complex transform with zero imaginary part}

The least complicated technique for computing a multidimensional
real/complex FFT is to simply compute a complex/complex FFT with zero
imaginary part, as in subsection~\ref{embed1d}.  The memory footprint
of this technique can be reduced by using the planar complex format.
This works in all cases, and may actually be fairly efficient for very
small problem sizes, particularly if the planar complex format is
available.

\subsection{Computation via packed 1D FFTs}

In subsection~\ref{packed1d}, we demonstrated a technique to compute
two real/complex 1D transforms using one complex/complex transform.
For multidimensional transforms, which start with multiple 1D
transforms to perform in the last direction, we can effectively batch
two 1D real/complex transforms together.  If the total problem size is
odd, there will be one transform leftover which can be computed by
embedding the real data into a complex array as in
subsection~\ref{embed1d}.  If the other dimensions are not very small,
then this shouldn't be too costly as compared with the even case.

The planar data format is important for this case, as it allows one to
perform complex/complex transforms in the last dimension of a
multidimensional without having to copy the data to a separate buffer.

\subsection{Computation for even-length data}

The algorithms for the real-complex transforms given in
subsection~\ref{even1d} can be used in the multi-dimensional case if
the last dimension is even.

For multi-dimensional transforms, we need to apply this in the last
dimension multiple times.  For batched transforms, this might be
annoying, as the 1D transform is already batched when applied to a 2D
(or 3D) transform, and then batching the 2D (or 3D) transform will
likely introduce a non-uniform distance parameter, though is appears
that we may have a solution using thread blocks.  Another possibility
is to create multiple pre/post-kernel nodes in the tree.

For a complex-complex 2D,
$\Four=\Four_x\circ\Four_y=\Four_y\circ\Four_x$, ie the sub-transforms
commute.  This isn't the case for a 2D real-complex transform when
taking advantage of Hermitian symmetry for the output.  To see this,
note that the input of a 2D real-to-complex FFT has $N_x\times{}N_y$
real values, the output is $N_x\times{}\(\floor{N_y/2} + 1\)$ complex
values, which is not symmetric in $N_x\leftrightarrow{}N_y$.  The
real-to-complex transform must then consist of first performing $N_x$
real-to-complex in the $y$ direction, followed by $\floor{N_y/2} + 1$
complex-to-complex transforms in the $x$ direction.
In other words, since
\begin{dmath}\label{nocommute}
  \Four_{x}\circ\text{post-process}_y\circ\Four_{y}
  \neq
  \Four_{x}\circ\Four_{y}\circ\text{post-process}_y
\end{dmath}
we cannot use the fact that $\Four_{x}\circ\Four_{y} = \Four_{xy}$.
It is easy to see why equation~\eqref{nocommute} is true, since
\begin{align}
  \Four_{y}
  &: N_x\times\frac{N_y}{2}\rightarrow N_x\times \frac{N_y}{2}\\
  \text{post-process}
  &: N_x\times\frac{N_y}{2}\rightarrow N_x\times \(\frac{N_y}{2}+1\)\\
  \Four_{x}
  &: N_x\times\(\frac{N_y}{2}+1\)\rightarrow N_x\times \(\frac{N_y}{2} +1\)
\end{align}
so, for the dimensions to match, we must perform $\Four_{y}$, then
$\text{post-process}$, and then $\Four_{x}$.

For a 2D complex-to-real transform, one performs the inverse
complex-to-real transform in the $y$ direction $N_x$ times, then
performs $N_y/2$ inverse complex-to-complex transforms in the $x$
direction.  In inverse case, however, one can make use of a 2D
complex-complex transform: the pre-processing can be applied, taking
the data from size $N_x\times{}\(Ny/2+1\)$ to $N_x\times{}Ny/2$, which
can be then transformed directly using a 2D complex-to-complex inverse
FFT.  This makes use of the fact that
\begin{dmath}
  \Four_{x}^{-1}  \circ  \Four_{y}^{-1} \circ \text{pre-process}_y
  = \Four_{xy}^{-1}\circ\text{pre-process}_y
  = \Four_{yx}^{-1}\circ\text{pre-process}_y
\end{dmath}
where the Fourier transforms are complex-complex.  Looking at
\begin{align}
  \text{pre-process}_y
  &: N_x\times\(\frac{N_y}{2}+1\)\rightarrow N_x\times \frac{N_y}{2}\\
  \Four_{y}^{-1}
  &: N_x\times\frac{N_y}{2}\rightarrow N_x\times \frac{N_y}{2}\\
  \Four_{x}^{-1}
  &: N_x\times\frac{N_y}{2}\rightarrow N_x\times \frac{N_y}{2}
\end{align}
one can see that, once the pre-processing is applied, all of the
dimensions match.

For 3D and higher transforms of dimension $d$, the direct transform
algorithm is
\begin{dmath}
  \Four_{d} \circ \text{post-process}_d
  \circ \Four_{d-1} \circ \dots \circ \Four_1
  =
  \Four_{d} \circ \text{post-process}_d  \circ \Four_{1 \dots (d-1)},
\end{dmath}
ie one performs $N_1\times{}N_{d-1}$ FFTs and post-process in the
dimension $d$, followed by a $N_d$ $d-1$-dimensional transforms in the
other dimensions.  For the inverse transform, we have
\begin{dmath}
  \text{pre-process}_d \circ \Four_d \circ \dots \circ \Four_1
  =
  \text{pre-process}_d \circ \Four_{d\dots{}1}
\end{dmath}
which is simply pre-processing in dimension $d$ followed by a
complex-complex transform on the entire array.

\end{document}

%%  LocalWords:  FFTs rocFFT FFT iy iY packedrcfft packinvr packinvNr
%%  LocalWords:  iO Hermitian ie rcfftpost XN pre FFTW cuFFT rocfft
%%  LocalWords:  Bluestein Multi multi prepostPlan xy nocommute Ny yx

